\scanpage{266}
- \mspage{84} -
\margin{IX}
\uline{Proposition} 7.1.5. \uline{Avec les notations précédentes pour} \(S,A,\hat B,\) \uline{supposons que} \(S\) \uline{soit} normal. \uline{Alors} \(\hat B\) \uline{est algébrisable, et le schéma abélien} \(B\) \uline{sur} \(S\) \uline{dont il provient est projectif sur} \(S.\)

On peut évidemment supposer \(A=A^0,\) et \(S\) intègre. D'après M. RAYNAUD [23 XI 1.13], l'hypothèse que \(S\) est normal implique que \(A\) est quasi-projectif sur \(S.\) Soit donc \(\uline{L}\) un Module inversible sur \(A\) qui est relativement ample, et formons le Module birigidifié
\(\delta(\uline{L})=\pi^*(\uline{L})\operatorname{pr}_1^*(\uline{L})^{-1}\operatorname{pr}_2^*(\uline{L})^{-1}\)
sur \(A\times_S A\) comme dans VIII 4.12. Comme les fibres de \(A_0\) sont de rang unipotent nul, il en est de même de celles de \(A\) (résulte facilement de 3.1), donc en vertu de VIII 7.5, \(\delta(\uline{L})\) définit une biextension \(W\) de \((A,A),\) évidemment symétrique, qui en vertu de VIII 4.13 (i.e. essentiellement en vertu de RAYNAUD [23 XI 1.11]) satisfait à la condition d'amplitude voulue en les points de \(S_0.\)

\uline{Remarque} 7.1.6. On peut prouver, sans hypothèse de normalité sur \(S,\) que \(\hat B\) est toujours algébrisable, en se ramenant au cas normal comme dans 13.5.3 a).

7.2. Notons maintenant le

\uline{Lemme} 7.2.1. \uline{Soient} \(S\) \uline{le spectre d'un anneau noethérien} \(V,\) \uline{séparé et complet pour la topologie définie par un idéal} \(\uline{m}\) \uline{de} \(V,\) \(C\) \uline{la catégorie des schémas en groupes} \(A\) \uline{sur} \(S,\) \uline{lisses de type fini sur} \(S,\) \uline{tels que} \(A/A^0\) \uline{soit un schéma en groupes étale fini sur} \(S\) \uline{et que} \(A^0\) \uline{soit une extension d'un schéma abélien} \(B\) \uline{par un tore isotrivial} \(T\) (SGA 3 IX 1.1),
\margin{402}
\margin{265}
\margin{402}
